\uuid{0wpC}
\exo7id{455}
\titre{exo7 455}
\auteur{cousquer}
\organisation{exo7}
\datecreate{2003-10-01}
\isIndication{false}
\isCorrection{false}
\chapitre{Propriétés de R}
\sousChapitre{Les rationnels}
\cpge{mpsi-mp2i.nombres-reels-suites.a, pcsi.raisonnement-vocabulaire-ensembliste.c, ptsi.raisonnement-vocabulaire-ensembliste.c, mpsi-mp2i.nombres-reels-suites.d, pcsi.nombres-reels-suites.c, ptsi.nombres-reels-suites.c, mpsi-mp2i.nombres-reels-suites.g, bcpst1.nombres-reels, bcpst1.suites-reelles}
\module{Analyse}
\niveau{L1}
\difficulte{}

\contenu{
\question{
On considère les nombres rationnels inférieurs à $\sqrt{2}$.
Y a-t-il un nombre rationnel juste avant $\sqrt{2}$, plus grand que tous les
nombres rationnels inférieurs à $\sqrt{2}$ ?

\noindent Une suite de nombres rationnels a-t-elle pour limite un nombre rationnel ?

\noindent Une suite de nombres décimaux a-t-elle pour limite un nombre décimal ?
}
}
